\chapter{The Mystery of Consciousness}

\section{Alice as Informational Structure}

We begin from the only thing that cannot be denied: Alice exists.

Based on the \emph{Four Initial Assumptions} \cite{meskanen2008,}, Alice must be an axiomatic system.

Let us digitize Alice’s DNA and simulate it on a computer. Alice is generated by an explicit algorithm: a DNA-like simulation code. A virtual Alice appears in a virtual universe. Due to substrate invariance consciousness arises during execution.

We begin a transformation by progressively replacing computation with lookup tables: less “code,” more “data.” At every step, nothing about Alice is allowed to change—not her memories, not her thoughts, not her pain. Equivalence within
axiomatic systems is absolute.

We continue this process to its limit. The program is fully unrolled. The code vanishes. What remains is a static data structure encoding the entire execution trace.
There is nothing left for the computer to run. If the simulation was multi-threaded, the static trace is highly fragmented, and Alice's bits are scattered within massive irrelevant data that looks like pure noise. 

Now the critical question arises:

\emph{Is Alice still conscious? If so, how does she “find” the bits that describe her?}

The answer is simple—and unavoidable:

\textbf{She doesn’t pick them. She is them.}

Alice is fully present within the relational structure of the trace. Time is not something that happens to her; it is encoded as the ordering of her states.
Memory, anticipation, and awareness all exist relationally within this structure. From the outside, it is static bits. From the inside, it is experience itself.

If consciousness could vanish at some ratio of code to data, that ratio would instantly become the most famous physical constant: the code/data threshold for vanishing pain.

Therefore, assuming the four initial assumptions hold, consciousness cannot depend on procedural execution.

This forces the conclusion:

\textbf{Consciousness is a property of informational structure, not of algorithmic execution.}

The usual objection arises: “But nothing is happening if the computer isn’t running.” This assumes time is fundamental. It is not. The trace encodes the order; experience is entirely relational. Nothing more is required.

Code is ultimately just data—bits, nothing more. If we could someday manage to extract the "pain function" from the DNA simulation code, it would amount to saying: this particular arrangement of bits “feels” pain, and this one does not.

It would be like saying that of two paintings on a wall, one is enduring unbearable agony and the other is happy.

In addition to software developers avoiding accidentally writing pain-aware software, artists should also be careful not to paint pain!



\section{Neural Networks}

Modern neural networks make this idea concrete. What appears as dynamic behavior is stored in a compressed, largely static structure.

The “computation” is not a step-by-step simulation, but the unfolding of constraints already embedded in the data.

Structure replaces procedure. Compression replaces simulation. And yet, even in static form, complex behavior, reasoning, and prediction emerge.

The distinction between “running” and “not running” becomes secondary.


\section{Free Will Restored}

Earlier, we wondered: could software ever allow free will? Deterministic code—branching if-else statements, fixed CPU instructions—appears to preclude choice.
The computer just runs the software; it cannot “choose otherwise,” regardless of what it might feel like to be Alice executing that program.

However, as we have demonstrated, consciousness can arise from static noisy information.

Consider a simple line of code:

\begin{verbatim}
if (a < 3){
    A
} else {
    B
}
\end{verbatim}

Deterministic? Yes. Sequential? Yes. But what does `3` really mean? In two’s complement, it might be interpreted as `3` or `-1`. As a floating-point number, it could be entirely different.  

Moreover, the code itself is nothing but data. Programmers interpret it as “instructions,” just as integers can be signed or unsigned.
There is no external device that determines the order in which bits get read. All the $2^n$ arrangements are equally real. 

As the number of bits \(n\) increases, the number of consistent interpretations grows combinatorially. Each interpretation defines a slightly different “Alice” inside the informational structure.
Some of these preserve consciousness; some do not.  Some may choose $A$, some $B$.

Alice is not forced to flow through this code like a train on tracks. All possible versions already exist, corresponding to every consistent interpretation of the bits.
Execution does not create Alice; it merely reveals one perspective among many.  

Which version does Alice experience? Most likely, the one corresponding to the most compressible, smooth, and coherent description—the version that dominates the observer-conditioned measure:

\[
\mathbb{P}(\text{Alice}_i) \propto 2^{-\mathcal{C}[\text{Alice}_i]}.
\]

Determinism at the CPU level does not constrain subjective experience. Execution is selection, not creation. And free will and pain? 

Alice navigates the space of all consistent interpretations, experiencing deliberation, anticipation, pain, and choice—not because the software code grants it,
but because the most compressible configuration.

Reasoning, and pain, is what it feels to be the most compressible, and hence, most probable Alice.


\section{Conclusion}

\begin{itemize}
  \item Alice exists
  \item Alice can be simulated
  \item The simulation can be transformed into static data
  \item Alice is preserved under that transformation
  \item Therefore, execution is irrelevant
  \item Therefore, time is internal to Alice
  \item Therefore, reality is structure rather than code
  \item The most compressible descriptions of Alice dominate the measure
  \item Threfore, Alice finds herself from the a law-like universe fine tuned to support her existence
  \item Because predictable information compresses best
\end{itemize}

Predictability is what we call the laws of physics.

